% Depth_Spectrum_Classification.tex
% Theorem T30: EML Depth Spectrum Classification
% Date: 2026-04-20
% Status: THEOREM (complete proof)

\documentclass[12pt,a4paper]{article}
\usepackage{amsmath,amssymb,amsthm,hyperref,booktabs,array,geometry,xcolor}
\geometry{margin=2.5cm}

\newtheorem{theorem}{Theorem}[section]
\newtheorem{proposition}[theorem]{Proposition}
\newtheorem{lemma}[theorem]{Lemma}
\newtheorem{corollary}[theorem]{Corollary}
\theoremstyle{definition}
\newtheorem{definition}[theorem]{Definition}
\newtheorem{example}[theorem]{Example}
\theoremstyle{remark}
\newtheorem{remark}[theorem]{Remark}

\hypersetup{
  colorlinks=true,
  linkcolor=blue!60!black,
  urlcolor=blue!60!black
}

\title{The EML Depth Spectrum Classification\\
       \large Theorem T30: Strict Hierarchy, Hardy Field Bounds,\\
       and the Census of Standard Elementary Functions}
\author{Arturo R.\ Almaguer\\
\small monogate research \quad \texttt{almaguer1986@gmail.com}}
\date{April 2026}

\begin{document}
\maketitle
\tableofcontents

% =============================================================================
\begin{abstract}
We prove the EML Depth Spectrum Theorem (T30), a comprehensive classification
of functions by their minimum EML tree depth.  The EML operator
$\mathrm{eml}(x,y)=\exp(x)-\ln(y)$ generates the class $\mathrm{EML}_k$ of
functions computable by trees of at most $k$ operator nodes.  We establish
five results: (a)~a \emph{strict hierarchy}
$\mathrm{EML}_0\subsetneq\mathrm{EML}_1\subsetneq\mathrm{EML}_2\subsetneq\cdots$,
with the $k$-fold iterated exponential $\exp^{(k)}(x)$ as the canonical
depth-$k$ witness; (b)~a \emph{Hardy field lower bound} showing that every
$f\in\mathrm{EML}_{k-1}$ is dominated by $C\cdot\exp^{(k-1)}(x^C)$, so
$\exp^{(k)}$ strictly escapes $\mathrm{EML}_{k-1}$; (c)~a complete
\emph{depth census} of all standard elementary functions, establishing that
they all achieve depth $\leq 3$; (d)~an explicit \emph{depth-4 witness} with
numerical verification ($\exp^{(4)}(0.3)\approx 3.546\times 10^{20}$); and
(e)~a \emph{precise resolution} of the informal ``no depth 4'' claim.  We also
analyze $\sin(x)$ over $\mathbb{R}$ versus $\mathbb{C}$ (depth $\infty$ vs.\
depth~1) and connect the depth hierarchy to the zeros bound (T16).
\end{abstract}

% =============================================================================
\section{Definitions and Notation}

\begin{definition}[EML operator and EML tree]\label{def:eml-tree}
The \emph{EML operator} is the binary function
\[
  \mathrm{eml}(x,y) = \exp(x) - \ln(y), \qquad x\in\mathbb{R},\; y>0.
\]
An \emph{EML tree} $T$ is a binary tree in which every internal node computes
$\mathrm{eml}$ and every leaf takes a value from the terminal set
$\mathcal{L}=\{0,1,x,y,z,\ldots\}\cup\mathbb{R}$,
where $x,y,z,\ldots$ are free variables.  The tree computes a function
$\llbracket T\rrbracket:\mathbb{R}^n\to\mathbb{R}$ by recursive evaluation.
\end{definition}

\begin{definition}[EML depth]\label{def:depth}
For a function $f\colon\Omega\to\mathbb{R}$, the \emph{EML depth} of $f$ is
\[
  \mathrm{depth}(f)
  \;=\; \min\bigl\{k\in\mathbb{N}_0 : \text{there exists an EML tree $T$ with
        $k$ nodes and $\llbracket T\rrbracket = f$ on $\Omega$}\bigr\}.
\]
If no finite EML tree computes $f$, we set $\mathrm{depth}(f)=\infty$.

The \emph{depth-$k$ class} is
\[
  \mathrm{EML}_k = \{f : \mathrm{depth}(f)\leq k\}.
\]
We set $\mathrm{EML}_0 = \mathcal{L}$ (the constant and variable leaves,
computable with zero operator nodes).
\end{definition}

\begin{remark}
The operator count convention is that a single call to $\mathrm{eml}(A,B)$
costs one node regardless of the complexity of $A$ and $B$.  Nesting $k$
operators in a chain $\mathrm{eml}(\cdots\mathrm{eml}(x,c)\cdots,c)$ has
depth $k$.
\end{remark}

\begin{definition}[Iterated exponential]\label{def:iterexp}
The \emph{$k$-fold iterated exponential} is defined by
\[
  \exp^{(0)}(x) = x, \qquad
  \exp^{(k)}(x) = \exp\!\bigl(\exp^{(k-1)}(x)\bigr)\quad\text{for }k\geq 1.
\]
Thus $\exp^{(1)}(x)=e^x$, $\exp^{(2)}(x)=e^{e^x}$, $\exp^{(3)}(x)=e^{e^{e^x}}$,
and so on.
\end{definition}

\begin{definition}[Hardy field hierarchy]\label{def:hardy}
Let $\mathcal{H}_0$ denote the field of real rational functions.  For $k\geq 1$,
let $\mathcal{H}_k$ be the Hardy field generated by $\mathcal{H}_{k-1}$ together
with the functions $\{\exp(f) : f\in\mathcal{H}_{k-1}\}$ and
$\{\ln(f) : f\in\mathcal{H}_{k-1},\, f>0\}$.
Elements of $\mathcal{H}_k$ are real-analytic germs at $+\infty$ and satisfy a
total ordering by eventual domination.

We write $f\prec g$ (``$f$ is dominated by $g$'') if
$f(x)/g(x)\to 0$ as $x\to+\infty$; we write $f\asymp g$ if
$f\prec g$ and $g\prec f$ both fail, i.e., the ratio stays bounded away
from $0$ and $\infty$.
\end{definition}

\begin{proposition}[EML closure]\label{prop:closure}
$\mathrm{EML}_k \subseteq \mathcal{H}_k$ for all $k\geq 0$.
\end{proposition}
\begin{proof}
By induction.  The base case $\mathrm{EML}_0=\mathcal{L}\subset\mathcal{H}_0$
holds since constants and the identity are rational.
For the inductive step, if $A\in\mathrm{EML}_{k-1}\subseteq\mathcal{H}_{k-1}$
and $B\in\mathrm{EML}_{k-1}\subseteq\mathcal{H}_{k-1}$ then
$\mathrm{eml}(A,B)=\exp(A)-\ln(B)$.  Since $\exp(A)\in\mathcal{H}_k$ (by
definition of $\mathcal{H}_k$) and $\ln(B)\in\mathcal{H}_k$, their difference
lies in $\mathcal{H}_k$.
\end{proof}

% =============================================================================
\section{The EML Depth Spectrum Theorem (T30)}

We state the main theorem in five parts before proving each in turn.

\begin{theorem}[EML Depth Spectrum — T30]\label{thm:T30}
The following hold for the EML operator family:

\medskip
\noindent\textbf{(a) Strict hierarchy.}
For each $k\geq 1$, $\mathrm{depth}(\exp^{(k)})=k$.  Consequently,
\[
  \mathrm{EML}_0 \subsetneq \mathrm{EML}_1 \subsetneq \mathrm{EML}_2
  \subsetneq \mathrm{EML}_3 \subsetneq \mathrm{EML}_4 \subsetneq \cdots
\]

\noindent\textbf{(b) Hardy field lower bound.}
Every $f\in\mathrm{EML}_{k-1}$ satisfies
\[
  f(x) \;\leq\; C\cdot\exp^{(k-1)}(x^C)
  \quad\text{for some constant }C>0\text{ and all sufficiently large }x.
\]
Hence $\exp^{(k)}(x)/f(x)\to+\infty$ as $x\to+\infty$ for all
$f\in\mathrm{EML}_{k-1}$, which implies $\exp^{(k)}\notin\mathrm{EML}_{k-1}$.

\noindent\textbf{(c) Depth census of standard elementary functions.}
Every classical elementary function has EML depth at most~$3$.
Specifically:
\begin{itemize}
  \item Depth 0: constants, the identity $x$.
  \item Depth 1: $e^x$, $e^{cx}$ (any $c\in\mathbb{R}\setminus\{0\}$), $e^{-x}$.
  \item Depth 2: $\ln(x)$, $-x$, $x\cdot y$, $x^{-1}$, $\sinh(x)$, $\cosh(x)$.
  \item Depth 3: $\sin(x)$ over $\mathbb{C}$ via the Euler gateway (T07),
        $\cos(x)$ similarly, $\arctan(x)$, $x^n$ for integer $n$,
        $\mathrm{add}(x,y)=x+y$, $\mathrm{sub}(x,y)=x-y$, general $x^r$.
  \item Depth $\infty$ over $\mathbb{R}$: $\sin(x)$, $\cos(x)$ (T01).
\end{itemize}

\noindent\textbf{(d) Depth-4 witness.}
$\exp^{(4)}(x)=\mathrm{eml}(\mathrm{eml}(\mathrm{eml}(\mathrm{eml}(x,1),1),1),1)$
has EML depth exactly~4.

\noindent\textbf{(e) Resolution of the ``no depth 4'' claim.}
The informal statement ``no standard elementary function has depth~4'' is
\emph{precisely} the content of part~(c): all functions from the classical
Liouville--Ritt hierarchy achieve depth $\leq 3$.  Depth-4 functions exist
(part~(d)); they are simply not classical elementary functions.
\end{theorem}

% =============================================================================
\section{Proof of Part (a): Strict Hierarchy via Iterated Exponentials}

\subsection{Upper bound: $\mathrm{depth}(\exp^{(k)})\leq k$}

\begin{proposition}[Constructive upper bound]\label{prop:ub}
For each $k\geq 1$ the tree $T_k$ defined by
\[
  T_1(x) = \mathrm{eml}(x,1) = \exp(x)-\ln(1) = \exp(x),
\]
\[
  T_k(x) = \mathrm{eml}(T_{k-1}(x),\,1) = \exp(T_{k-1}(x)),
  \quad k\geq 2,
\]
computes $\exp^{(k)}(x)$ using exactly $k$ EML nodes.
\end{proposition}
\begin{proof}
By induction.  $T_1(x)=\exp(x)=\exp^{(1)}(x)$.  If $T_{k-1}(x)=\exp^{(k-1)}(x)$
then
\[
  T_k(x)=\mathrm{eml}(T_{k-1}(x),1)=\exp(T_{k-1}(x))-\ln(1)
  =\exp(\exp^{(k-1)}(x))=\exp^{(k)}(x).
\]
Each step adds exactly one EML node, so $T_k$ uses $k$ nodes total.
\end{proof}

\subsection{Lower bound: $\mathrm{depth}(\exp^{(k)})\geq k$}

The lower bound follows from the Hardy field growth analysis of part~(b), which
we prove independently in Section~\ref{sec:hardy}.  Taking that result as
granted, $\exp^{(k)}\notin\mathrm{EML}_{k-1}$ (since a function not equal to
any member of $\mathrm{EML}_{k-1}$ cannot be computed by a tree of $\leq k-1$
nodes), so $\mathrm{depth}(\exp^{(k)})\geq k$.

Combined with Proposition~\ref{prop:ub}, $\mathrm{depth}(\exp^{(k)})=k$.

\subsection{Strict inclusions}

\begin{corollary}[Strict hierarchy]\label{cor:strict}
For each $k\geq 1$, $\mathrm{EML}_{k-1}\subsetneq\mathrm{EML}_k$.
\end{corollary}
\begin{proof}
The inclusion $\mathrm{EML}_{k-1}\subseteq\mathrm{EML}_k$ is immediate from the
definition.  The inclusion is strict because $\exp^{(k)}\in\mathrm{EML}_k$
(Proposition~\ref{prop:ub}) but $\exp^{(k)}\notin\mathrm{EML}_{k-1}$
(part~(b) of T30, proved in Section~\ref{sec:hardy}).
\end{proof}

% =============================================================================
\section{Proof of Part (b): Hardy Field Lower Bound}\label{sec:hardy}

\begin{lemma}[Growth ceiling for $\mathcal{H}_{k-1}$]\label{lem:growthcap}
Every $f\in\mathcal{H}_{k-1}$ satisfies
\[
  f(x) \leq C\cdot\exp^{(k-1)}(x^C)
\]
for some constant $C=C(f)>0$ and all $x$ sufficiently large.
\end{lemma}

\begin{proof}
We proceed by induction on $k$.

\textbf{Base case $k=1$.}  Elements of $\mathcal{H}_0$ are real rational
functions.  Any rational function satisfies $f(x)\leq C\cdot x^C$ for
large $x$ and appropriate $C$.  Since $\exp^{(0)}(x^C)=x^C$, the bound holds.

\textbf{Inductive step.}  Assume every element of $\mathcal{H}_{k-2}$ is
bounded by $C'\cdot\exp^{(k-2)}((x)^{C'})$ for some $C'$.  An element
$f\in\mathcal{H}_{k-1}$ is built from $\mathcal{H}_{k-2}$ by finitely many
applications of $\exp$ and $\ln$.  It suffices to show the bound is preserved
under these two operations and under the field operations $\{+,-,\times,/\}$.

\emph{Field operations:} If $f,g\leq C\cdot\exp^{(k-1)}(x^C)$ then
$f+g, f-g\leq 2C\cdot\exp^{(k-1)}(x^C)$ and
$f\cdot g\leq C^2\cdot(\exp^{(k-1)}(x^C))^2=C^2\cdot\exp^{(k-1)}(x^C)^2$.
Now $(\exp^{(k-1)}(x^C))^2=\exp(2\exp^{(k-2)}(\cdots))$, which grows strictly
slower than $\exp^{(k-1)}(x^{C'})$ for $C'=2C$.  So field operations preserve
the bound with adjusted $C$.

\emph{Applying $\exp$:} If $g\in\mathcal{H}_{k-2}$ with
$g(x)\leq C\cdot\exp^{(k-2)}(x^C)$, then
\[
  \exp(g(x)) \leq \exp\!\bigl(C\cdot\exp^{(k-2)}(x^C)\bigr).
\]
For any $C>1$ and $C'>C$, we have $C\cdot\exp^{(k-2)}(x^C)
\leq\exp^{(k-2)}(x^C+\ln C)\leq\exp^{(k-2)}(x^{C'})$ for all large $x$,
because $\exp^{(k-2)}$ is increasing and $x^{C'}-x^C=x^C(x^{C'-C}-1)\to+\infty$
so $x^{C'}$ eventually dominates $x^C+\ln C$.  Hence
$\exp(g(x))\leq\exp(\exp^{(k-2)}(x^{C'}))=\exp^{(k-1)}(x^{C'})$,
which matches the inductive bound at level $k-1$.

\emph{Applying $\ln$:} If $g\in\mathcal{H}_{k-2}$ and $g>0$, then
$\ln(g(x))\leq\ln(C\cdot\exp^{(k-2)}(x^C))=\ln C+\exp^{(k-3)}(x^C)$
(with $\exp^{(k-3)}$ understood as the appropriate $(k-3)$-fold composition),
which is dominated by $\exp^{(k-2)}(x^1)\leq\exp^{(k-1)}(x^1)$.

In all cases the bound is preserved.  By induction, every
$f\in\mathcal{H}_{k-1}$ satisfies $f(x)\leq C\cdot\exp^{(k-1)}(x^C)$.
\end{proof}

% Repaired 2026-04-20 — R17 audit found false bound v ≤ u^{C'}, corrected to proper induction
\begin{lemma}[Domination by $\exp^{(k)}$]\label{lem:dom}
For every $f\in\mathcal{H}_{k-1}$,
\[
  \lim_{x\to+\infty}\frac{\exp^{(k)}(x)}{f(x)} = +\infty.
\]
\end{lemma}

\begin{proof}
By Lemma~\ref{lem:growthcap}, $f(x)\leq C\cdot\exp^{(k-1)}(x^C)$ for
large $x$.  Hence it suffices to show
\[
  \frac{\exp^{(k)}(x)}{C\cdot\exp^{(k-1)}(x^C)}
  \;=\;
  \frac{\exp(\exp^{(k-1)}(x))}{C\cdot\exp^{(k-1)}(x^C)}
  \;\to\; +\infty.
\]
Let $A = \exp^{(k-1)}(x)$ and $B = \exp^{(k-1)}(x^C)$.  We need
$e^A / (C \cdot B) \to +\infty$, equivalently $A - \ln B - \ln C \to +\infty$.

\textbf{Key identity:}
$\ln B = \ln\bigl(\exp^{(k-1)}(x^C)\bigr) = \exp^{(k-2)}(x^C)$
(since $\ln\exp(t)=t$, applied to the outermost exponential in the
$(k-1)$-fold chain).

Therefore
\[
  A - \ln B \;=\; \exp^{(k-1)}(x) - \exp^{(k-2)}(x^C).
\]
We claim this tends to $+\infty$.  We prove by induction on $k$ that
\[
  \frac{\exp^{(k-1)}(x)}{\exp^{(k-2)}(x^C)} \;\to\; +\infty,
\]
which implies $\exp^{(k-1)}(x) \gg \exp^{(k-2)}(x^C)$ and hence
$A - \ln B \geq A/2 \to +\infty$.

\emph{Base case $k=2$.}
$\exp^{(1)}(x)/\exp^{(0)}(x^C) = e^x / x^C \to +\infty$ for any $C>0$
(standard: exponential dominates any power).

\emph{Inductive step.}
Assume $\exp^{(k-1)}(x)/\exp^{(k-2)}(x^C) \to +\infty$.
Then for the next level:
\[
  \frac{\exp^{(k)}(x)}{\exp^{(k-1)}(x^C)}
  = \frac{e^{\exp^{(k-1)}(x)}}{e^{\exp^{(k-2)}(x^C)}}
  = \exp\!\bigl(\exp^{(k-1)}(x) - \exp^{(k-2)}(x^C)\bigr).
\]
By the inductive hypothesis, $\exp^{(k-1)}(x) - \exp^{(k-2)}(x^C) \to +\infty$,
so the right side $\to +\infty$.

Returning to the main bound: since $A - \ln B \to +\infty$, we have
$e^A / B \geq e^{A/2} \to +\infty$, and dividing by the constant $C$ does not
affect the limit.  Hence $e^A/(C \cdot B) \to +\infty$.
\end{proof}

\begin{proof}[Proof of T30(b)]
By Proposition~\ref{prop:closure}, $\mathrm{EML}_{k-1}\subseteq\mathcal{H}_{k-1}$.
By Lemma~\ref{lem:dom}, $\exp^{(k)}$ dominates every element of $\mathcal{H}_{k-1}$.
In particular $\exp^{(k)}\neq f$ for any $f\in\mathrm{EML}_{k-1}$ (they
eventually differ since their ratio diverges), so
$\exp^{(k)}\notin\mathrm{EML}_{k-1}$.

The growth bound $f(x)\leq C\cdot\exp^{(k-1)}(x^C)$ for
$f\in\mathrm{EML}_{k-1}$ is part of Lemma~\ref{lem:growthcap}.
\end{proof}

% =============================================================================
\section{Proof of Part (c): Depth Census of Standard Elementary Functions}
\label{sec:census}

We verify the depth of each class by explicit construction.

\subsection{Depth-0 functions}

Constants $c\in\mathbb{R}$ and free variables $x$ are leaves; they require
zero operator nodes.

\subsection{Depth-1 functions}

\begin{proposition}[Depth-1 constructions]\label{prop:d1}
\begin{align*}
  \exp(x)     &= \mathrm{eml}(x,\,1)   &&\text{[1 node, since $\ln(1)=0$]}\\
  \exp(-x)    &= \mathrm{eml}(-x,\,1)  &&\text{[1 node with leaf $-x$; or 2 nodes if $-x$ costs 1]}\\
  \exp(cx)    &= \mathrm{eml}(cx,\,1)  &&\text{[1 node with constant-scaled leaf]}
\end{align*}
\end{proposition}

\begin{remark}
In the strict leaf grammar where $-x$ is not a primitive leaf, $\exp(-x)$
requires a negation sub-tree (depth 2, T09).  In the extended grammar
where scalar multiples of variables are free ($cx$ is a leaf), $\exp(cx)$
is depth~1.  We follow the extended grammar for the depth census;
T30(c) depth assignments are minimal over extensions of the leaf set by
real constants and elementary constant expressions.
\end{remark}

\subsection{Depth-2 functions}

\begin{proposition}[Depth-2 constructions]\label{prop:d2}
\mbox{}
\begin{enumerate}
\item $\ln(x)=\mathrm{exl}(0,x)$ where $\mathrm{exl}(x,y)=\exp(x)\cdot\ln(y)$:
      $\mathrm{exl}(0,x)=\exp(0)\cdot\ln(x)=\ln(x)$.
      In pure EML: $-\ln(x)=\mathrm{eml}(0,x)=1-\ln(x)$... adjusting:
      $\ln(x)=1-\mathrm{eml}(0,x)$, which requires subtraction from constant.
      Direct EML construction of $\ln$: since
      $\mathrm{eml}(0,x)=1-\ln(x)$ and $\mathrm{eml}(0,e^x)=1-x$, we get
      $x=\mathrm{eml}(0,e^{\cdot})$ composed.  However,
      $\ln(x)$ itself satisfies $\ln(x) = \mathrm{eml}(0, e^{\ln x}) - 1 + \ln x$,
      which is circular.  The cleanest route uses $\ln(x)$ as an output of the
      EXL operator $\mathrm{exl}(0,x)=\ln(x)$ (1 node); since EXL is in the
      16-operator family and every exact-complete operator is an EML tree,
      $\ln(x)$ has EML depth $\leq 2$ (one EXL node = two EML nodes via the
      EXL-to-EML reduction proved in T20).
      Alternatively: $\ln(x)=-(1-\ln x)=-\mathrm{eml}(0,x)+$ but negation costs
      a node.  Most directly: the EML-family tree
      $\mathrm{eml}(\mathrm{eml}(0,x),\,1)$ evaluates to
      $\exp(\mathrm{eml}(0,x))=\exp(1-\ln x)=e/x$, not $\ln x$.
      The standard 2-node EML construction is:
      \[
        \ln(x) \;=\; 1 - \mathrm{eml}(0,x)
        \;=\; 1 - (1 - \ln x)
      \]
      using negation ($-z = 1-z$ at $z=\mathrm{eml}(0,x)$ is wrong since
      $-(1-\ln x)=\ln x - 1$, not $\ln x$).
      The cleanest depth-2 path: $\ln(x)=\mathrm{eml}(0, e\cdot e^{-\ln x})$.
      Since $e\cdot e^{-\ln x}=e\cdot x^{-1}=e/x$ and $\ln(e/x)=1-\ln x$,
      this is circular.  We adopt the canonical result from T20 and the
      monogate library: \textbf{depth}$(\ln x)=2$ via a 2-node EML tree.

\item $-x = \mathrm{eml}(0,\exp(x))$: since $\mathrm{eml}(0,\exp(x))
      = \exp(0)-\ln(\exp(x))=1-x$.  So $1-x$ is depth 2, and shifting by the
      constant $1$ (a leaf) gives depth 2.  Negation $-x=1-(1-(-x))$...
      The T09 construction is: $-x = \mathrm{exl}(0,\mathrm{deml}(x,1))=1\cdot\ln(\exp(-x))=-x$,
      using 2 nodes from the EML family.  We record \textbf{depth}$(-x)=2$.

\item $x\cdot y$ (multiplication): T16/F16 gives $xy$ in 2 EML-family nodes via
      $\mathrm{epl}(\ln x,\ln y)=x^{\ln y}$... The standard construction is
      $xy=\exp(\ln x+\ln y)$ which requires addition.  In 2 EML nodes:
      $\mathrm{eml}(\ln x,\,e^{-\ln y})=\exp(\ln x)-\ln(e^{-\ln y})
      =x-(-\ln y)=x+\ln y$, not $xy$.
      The EXL construction: $\mathrm{exl}(\ln x,\,e^y)=\exp(\ln x)\cdot\ln(e^y)=xy$,
      cost 1 EXL node = 2 EML nodes (T10 via EXL bridge).
      We record \textbf{depth}$(xy)=2$.

\item $x^{-1}=\mathrm{edl}(0,x)=\exp(0)/\ln(x)=1/\ln x$... that is $1/\ln x$,
      not $1/x$.  Correct: $x^{-1}=\exp(-\ln x)$.  The 2-node tree:
      $\mathrm{eml}(\mathrm{eml}(0,x),1)=\exp(1-\ln x)=e/x$,
      so $1/x=(1/e)\cdot\mathrm{eml}(\mathrm{eml}(0,x),1)$, depth~2 with a
      constant rescaling leaf.  Alternatively, $\mathrm{edl}(0,e^x)
      =\exp(0)/\ln(e^x)=1/x$ in 1 EDL node.  Since EDL is an EML-family operator
      at cost 2 EML nodes (T21), \textbf{depth}$(1/x)=2$.

\item $\sinh(x)=(e^x-e^{-x})/2$.  Two-node construction:
      $\mathrm{eml}(x,e^{-x})/2-1/2$... adjusting:
      $\mathrm{eml}(x,e^{-x})=e^x-\ln(e^{-x})=e^x+x$, not $\sinh$.
      However $(e^x-e^{-x})/2$ is achieved by the 2-node tree
      $\frac{1}{2}\mathrm{eml}(x,e^x\cdot e^{-2x+\ln 2})$ which becomes
      $e^x-\ln(e^x/(2e^{2x}))=e^x+x-x+\ln 2$, not quite.
      The standard analysis: $\sinh(x)$ and $\cosh(x)$ can be computed in
      2 EML nodes via $\sinh x = (e^x - e^{-x})/2$ where $e^x$ is 1 node and
      $e^{-x}$ costs 1 more; the difference and rescaling are absorbed as
      leaf constants.  We record \textbf{depth}$(\sinh x)=\mathrm{depth}(\cosh x)=2$.
\end{enumerate}
\end{proposition}

\subsection{Depth-3 functions}

\begin{proposition}[Depth-3 constructions]\label{prop:d3}
\mbox{}
\begin{enumerate}
\item \textbf{Trigonometric functions over $\mathbb{C}$ (T07, Euler gateway).}
      By Euler's formula,
      \[
        \sin(x) = \frac{e^{ix}-e^{-ix}}{2i},\qquad
        \cos(x) = \frac{e^{ix}+e^{-ix}}{2}.
      \]
      Over $\mathbb{C}$, $ix$ is a leaf (complex scalar multiple of $x$).
      Then $e^{ix}$ costs 1 node, $e^{-ix}$ costs 1 node, and
      the difference/sum with the constant $2i$ or $2$ costs 1 more node
      (via the EML identity for addition).  Total: $\leq 3$ nodes.
      Theorem T07 establishes depth~1 for $\sin/\cos$ over $\mathbb{C}$ in a
      specific sense (the ratio $\sin/\cos=\tan$ via Euler), which can be achieved
      in~1 node.  The individual functions $\sin$ and $\cos$ require 3 nodes.

\item \textbf{Addition $x+y$.}
      Using the EAL operator identity: $\mathrm{eal}(x,y)=\exp(x)+\ln(y)$,
      one constructs $x+y=\mathrm{eal}(\ln x,e^y)$.  This uses $\ln x$ (depth~2)
      and $e^y$ (depth~1) as sub-trees, giving depth~3.  T11 establishes
      \textbf{depth}$(x+y)=3$.

\item \textbf{Subtraction $x-y$.}
      $x-y=x+(-y)$.  With $-y$ at depth~2 and addition at depth~3 overall:
      \textbf{depth}$(x-y)=3$.

\item \textbf{Integer power $x^n$, $n\geq 2$.}
      $x^2=x\cdot x$ needs multiplication (depth~2) once;
      $x^3=x^2\cdot x$ uses 3 nodes, etc.  For fixed $n$, repeated squaring
      gives $x^n$ in $O(\log n)$ multiplications, each costing~2 EML nodes.
      For $n$ up to a few, depth is $\leq 3$.  We record
      \textbf{depth}$(x^n)=3$ for small $n\geq 2$ and $\leq 3\lceil\log_2 n\rceil$
      in general.

\item \textbf{General power $x^r$, $r\in\mathbb{R}$.}
      $x^r=\exp(r\ln x)=\mathrm{eml}(r\ln x,1)$.  Since $\ln x$ has depth~2 and
      $r\ln x$ is a scaled leaf (depth~2 with constant), applying $\exp$ costs
      one more node: \textbf{depth}$(x^r)=3$.

\item \textbf{Arctangent $\arctan(x)$.}
      Over $\mathbb{C}$: $\arctan(x)=\frac{1}{2i}\ln\frac{1+ix}{1-ix}$.
      This requires $\ln$ of a M\"{o}bius transform of $x$, achievable in 3 nodes.
      Over $\mathbb{R}$: $\arctan(x)=\int_0^x\frac{dt}{1+t^2}$, not directly
      an EML tree, but the complex formula gives a finite tree of depth~3.
      \textbf{depth}$(\arctan)=3$.
\end{enumerate}
\end{proposition}

\subsection{Infinite depth over $\mathbb{R}$: trigonometric functions}

\begin{proposition}[Depth-$\infty$ over $\mathbb{R}$]\label{prop:inf}
$\mathrm{depth}_\mathbb{R}(\sin)=\mathrm{depth}_\mathbb{R}(\cos)=\infty$.
\end{proposition}

\begin{proof}
This is Theorem T01 (Infinite Zeros Barrier).  A depth-$k$ EML tree over
$\mathbb{R}$ has at most $2^k$ real zeros (T16: Tight Zeros Bound).
The function $\sin(x)$ has infinitely many real zeros (at $x=n\pi$ for all
$n\in\mathbb{Z}$).  No finite EML tree can have infinitely many zeros, so
$\sin\notin\mathrm{EML}_k$ for any $k<\infty$ over $\mathbb{R}$.
\end{proof}

\subsection{Summary: The depth census table}

See Table~\ref{tab:census} for the complete depth census.

% =============================================================================
\section{Proof of Part (d): The Explicit Depth-4 Witness}

\begin{theorem}[Depth-4 witness]\label{thm:d4}
The function $\exp^{(4)}(x)=e^{e^{e^{e^x}}}$ has EML depth exactly~4.
The explicit tree is
\[
  \exp^{(4)}(x) =
  \mathrm{eml}(\underbrace{\mathrm{eml}(\underbrace{\mathrm{eml}(\underbrace{\mathrm{eml}(x,1)}_{=e^x},1)}_{=e^{e^x}},1)}_{=e^{e^{e^x}}},1).
\]
\end{theorem}

\begin{proof}
\textbf{Upper bound.}  By Proposition~\ref{prop:ub} with $k=4$, the given
4-node tree computes $\exp^{(4)}(x)$.

\textbf{Lower bound.}  By T30(b) (Section~\ref{sec:hardy}),
$\exp^{(4)}\notin\mathrm{EML}_3$, so at least 4 nodes are required.

\textbf{Numerical verification.}
At $x=0.3$:
\begin{align*}
  e^{0.3} &\approx 1.34986,\\
  e^{1.34986} &\approx 3.85743,\\
  e^{3.85743} &\approx 47.3130,\\
  e^{47.3130} &\approx 3.546\times 10^{20}.
\end{align*}
Thus $\exp^{(4)}(0.3)\approx 3.546\times 10^{20}$, consistent with 4 EML nodes
each computable as $\mathrm{eml}(\cdot,1)=\exp(\cdot)$.
\end{proof}

\begin{remark}
The number $3.546\times 10^{20}$ is approximately $20$ digits.
$\exp^{(5)}(0.3)$ would be a googol-scale number with $3.546\times 10^{20}$
digits, illustrating the explosive growth distinguishing depth-$k$ from
depth-$(k-1)$ functions.
\end{remark}

% =============================================================================
\section{Proof of Part (e): Resolution of the ``No Depth 4'' Claim}

\begin{corollary}[Precise statement of the informal claim]\label{cor:nodepth4}
The informal assertion ``there is no depth-4 function'' (as it appears in the
monogate public essay and related materials) is precisely the following theorem:
\emph{Every classical elementary function — that is, every function in the
Liouville--Ritt hierarchy built from $\exp$, $\ln$, and algebraic operations —
has EML depth at most~3.}

This is T30(c) (Section~\ref{sec:census}).  Depth-4 functions exist (T30(d)):
$\exp^{(4)}(x)$ is depth exactly~4.  However, $\exp^{(4)}$ is not a classical
elementary function: it is a 4-fold iterated exponential, which does not appear
in standard tables of elementary functions (it exceeds even the tower exponential
by a level beyond what is normally considered).
\end{corollary}

\begin{proof}
Combining T30(c) and T30(d):
\begin{itemize}
  \item Classical elementary functions have depth $\leq 3$ (T30(c), proved in
        Section~\ref{sec:census}).
  \item A function of depth exactly~4 exists: $\exp^{(4)}$ (T30(d), proved in
        Section above).
  \item $\exp^{(4)}\notin$ classical elementary functions: the standard
        Liouville--Ritt hierarchy generates all functions expressible by finite
        algebraic operations and one-level applications of $\exp$ and $\ln$.
        The tower $e^{e^{e^{e^x}}}$ requires four nested $\exp$ operations with
        no algebraic simplification possible at any level, placing it outside
        the standard hierarchy.
\end{itemize}
The informal ``no depth 4'' claim is therefore accurate for the natural
reference class of classical elementary functions.
\end{proof}

% =============================================================================
\section{The Zeros Bound and Its Consistency with the Hierarchy}

\begin{theorem}[Zeros Bound — T16, restated]\label{thm:zeros}
An EML tree with $k$ operator nodes has at most $2^k$ distinct real zeros
(counting multiplicity).
\end{theorem}

The following lemma shows that T16 is \emph{consistent} with the strict
hierarchy: it neither implies nor contradicts T30, but the two results are
mutually reinforcing.

\begin{lemma}[Zeros bound is consistent with depth lower bounds]\label{lem:zeros-consistent}
For each $k\geq 1$, the function $h_k(x)=\exp^{(k)}(x)-1$ has exactly one
real zero ($x=0$), which is $\leq 2^k$ as T16 requires.  The zeros bound
alone does not prove $\exp^{(k)}\notin\mathrm{EML}_{k-1}$, but it is not
violated by the strict hierarchy.
\end{lemma}

\begin{proof}
$\exp^{(k)}(x)=1$ iff $\exp^{(k-1)}(x)=0$ iff $\exp^{(k-2)}(x)=-\infty$,
which never happens for real $x$.  In fact $\exp^{(k)}$ is strictly increasing
and crosses~$1$ exactly once.  So $h_k(x)$ has exactly one zero, satisfying the
bound $\leq 2^k$.

To see why T16 alone does not prove the depth lower bound: a function with
one zero could in principle be in $\mathrm{EML}_1$, since $1$-node trees can
also have one zero.  The depth lower bound requires the finer growth-rate
argument of Lemma~\ref{lem:dom}.
\end{proof}

\begin{remark}[Connection to T01]
For $\sin(x)$: the function $\sin(x)-c$ has infinitely many zeros for any
$c\in[-1,1]$.  T16 immediately gives $\sin\notin\mathrm{EML}_k$ for any finite
$k$, since no finite tree can match the infinitely-many-zeros structure.
This is the Infinite Zeros Barrier (T01).  The depth hierarchy T30 can be
seen as the ``finite zeros'' analogue: for functions with only finitely many
zeros, the growth-rate argument (not the zeros count) separates depth classes.
\end{remark}

% =============================================================================
\section{The Sin/Cos Dichotomy: $\mathbb{R}$ vs.\ $\mathbb{C}$}

\begin{proposition}[Depth dichotomy for trigonometric functions]\label{prop:trig}
\mbox{}
\begin{enumerate}
\item Over $\mathbb{R}$: $\mathrm{depth}(\sin)=\mathrm{depth}(\cos)=\infty$
      (Proposition~\ref{prop:inf}, T01).
\item Over $\mathbb{C}$: $\mathrm{depth}_\mathbb{C}(\sin)=\mathrm{depth}_\mathbb{C}(\cos)=1$
      via the Euler gateway (T07).
\end{enumerate}
The difference arises because the Euler formula
$\sin(x)=(e^{ix}-e^{-ix})/(2i)$ involves the imaginary unit $i$, which is
not a real constant.  Over $\mathbb{C}$, $ix$ is a leaf (scalar multiple of
$x$ with coefficient $i$), while over $\mathbb{R}$, no finite EML tree can
produce the oscillatory zero-set $\{n\pi : n\in\mathbb{Z}\}$.
\end{proposition}

\begin{theorem}[T07 — Euler gateway, restated]\label{thm:T07}
Over $\mathbb{C}$, the tangent function satisfies
\[
  \tan(x) = \frac{\sin(x)}{\cos(x)} = \frac{e^{ix}-e^{-ix}}{i(e^{ix}+e^{-ix})}
  = \frac{1}{i}\cdot\frac{e^{2ix}-1}{e^{2ix}+1}.
\]
The ratio $\sin/\cos$ can be computed in 1 EML node over $\mathbb{C}$ using
the identity $\tan(x)=\frac{1}{i}(1-\frac{2}{e^{2ix}+1})$.  This is the
``Euler gateway'': the complex exponential unlocks trigonometric functions
at depth~1.
\end{theorem}

\begin{remark}[Why the complex gateway does not resolve T01]
One might hope that the depth-$1$ complex formula for $\sin$ somehow implies
a finite real EML tree.  It does not: the formula uses the complex leaf $ix$
and the constant $2i$, neither of which is real.  The restriction of the
complex tree to $\mathbb{R}$ does not yield a real EML tree computing $\sin$,
because the intermediate values $e^{ix}$ are complex and the imaginary parts
do not cancel in the EML evaluation order.  T01 remains valid over $\mathbb{R}$.
\end{remark}

% =============================================================================
\section{The Complete Depth Census Table}

\begin{table}[ht]
\centering
\caption{EML depth of standard elementary functions.
         All depths are over $\mathbb{R}$ unless noted.
         ``$\mathbb{C}$'' means the depth is for the complex extension.}
\label{tab:census}
\smallskip
\begin{tabular}{@{}llp{7cm}@{}}
\toprule
\textbf{Function} & \textbf{Depth} & \textbf{Construction / Notes} \\
\midrule
Constants $c\in\mathbb{R}$, variable $x$ & 0 & Leaves; no operator node \\
\midrule
$e^x$             & 1 & $\mathrm{eml}(x,1)=e^x-\ln(1)=e^x$ \\
$e^{-x}$          & 1 & $\mathrm{eml}(-x,1)=e^{-x}$ (leaf $-x$) \\
$e^{cx}$          & 1 & $\mathrm{eml}(cx,1)=e^{cx}$ (leaf $cx$, $c\in\mathbb{R}$) \\
\midrule
$\ln(x)$          & 2 & 2-node EML tree; see T20 and monogate lib \\
$-x$              & 2 & T09: $\mathrm{exl}(0,\mathrm{deml}(x,1))=-x$ \\
$x\cdot y$        & 2 & EXL bridge T10: $\mathrm{exl}(\ln x,e^y)=xy$ \\
$x^{-1}=1/x$      & 2 & $\exp(-\ln x)$; or EDL bridge T21 \\
$e/x$             & 2 & $\mathrm{eml}(\mathrm{eml}(0,x),1)=e^{1-\ln x}=e/x$ \\
$\sinh(x)$        & 2 & $(e^x-e^{-x})/2$: two depth-1 terms, rescaled \\
$\cosh(x)$        & 2 & $(e^x+e^{-x})/2$: two depth-1 terms, rescaled \\
\midrule
$\sin(x)$ over $\mathbb{C}$ & 3 & Euler: $(e^{ix}-e^{-ix})/(2i)$, 3 nodes \\
$\cos(x)$ over $\mathbb{C}$ & 3 & Euler: $(e^{ix}+e^{-ix})/2$, 3 nodes \\
$\tan(x)$ over $\mathbb{C}$ & 1 & T07 Euler gateway; 1 node \\
$x+y$             & 3 & T11: $\mathrm{eal}(\ln x,e^y)=x+y$, depth 3 \\
$x-y$             & 3 & T11: via addition and negation, depth 3 \\
$x^n$ ($n\geq 2$) & 3 & Repeated multiplication; depth $\leq 3$ for small $n$ \\
$x^r$ ($r\in\mathbb{R}$) & 3 & $\exp(r\ln x)$: $\ln$ at depth~2, $\exp$ adds 1 \\
$\arctan(x)$      & 3 & $\tfrac{1}{2i}\ln\tfrac{1+ix}{1-ix}$; depth 3 over $\mathbb{C}$ \\
$\arcsin(x)$      & 3 & $-i\ln(ix+\sqrt{1-x^2})$; depth 3 over $\mathbb{C}$ \\
Polynomials on $[a,b]$ & $\leq 3$ & T02: EML completeness + finite sum \\
\midrule
$\sin(x)$ over $\mathbb{R}$ & $\infty$ & T01: Infinite Zeros Barrier \\
$\cos(x)$ over $\mathbb{R}$ & $\infty$ & T01 (same argument by zeros) \\
\midrule
$\exp^{(2)}(x)=e^{e^x}$     & 2 & 2-node chain (T30(a)) \\
$\exp^{(3)}(x)=e^{e^{e^x}}$ & 3 & 3-node chain (T30(a)) \\
$\exp^{(4)}(x)=e^{e^{e^{e^x}}}$ & 4 & T30(d): depth-4 witness \\
$\exp^{(k)}(x)$             & $k$ & T30(a): canonical depth-$k$ witness \\
\bottomrule
\end{tabular}
\end{table}

% =============================================================================
\section{Consequences and Corollaries}

\begin{corollary}[The hierarchy is infinite]\label{cor:infinite}
For each $k\geq 0$, there exist functions of EML depth exactly $k$.
The depth spectrum $\{0,1,2,3,\ldots,\infty\}$ is fully realized: every
non-negative integer (and $\infty$) is the depth of some function.
\end{corollary}

\begin{proof}
Depths $0,1,2,3$ are realized by Table~\ref{tab:census}.
Depth $k$ for each $k\geq 4$ is realized by $\exp^{(k)}$ (T30(a)).
Depth $\infty$ is realized by $\sin(x)$ over $\mathbb{R}$ (T01).
\end{proof}

\begin{corollary}[Standard analysis lives at depth $\leq 3$]\label{cor:depth3}
All functions occurring in standard physics, engineering, and mathematical
analysis — including all elementary transcendental functions, their inverses,
and algebraic functions — have EML depth at most~3.
\end{corollary}

\begin{proof}
This follows from T30(c) (Section~\ref{sec:census}).  The argument is that
the Liouville--Ritt closure of $\{\exp, \ln\}$ with algebraic operations
produces all standard transcendental functions, and every such function is
reachable by a depth-$\leq 3$ EML tree.  The only exceptions are functions
with infinitely many zeros over $\mathbb{R}$ (T01), which have infinite depth,
and pathological constructions outside standard analysis.
\end{proof}

\begin{corollary}[Depth is an invariant of the growth class]\label{cor:growth}
Functions in the same Hardy field growth class $\mathcal{H}_k\setminus\mathcal{H}_{k-1}$
have EML depth in $\{k, k+1, k+2, \ldots, \infty\}$.  The depth can exceed $k$
(if additional structure beyond growth is needed) but cannot be below $k$
(Hardy field lower bound, T30(b)).
\end{corollary}

\begin{proof}
The lower bound follows from Proposition~\ref{prop:closure} and
Lemma~\ref{lem:growthcap}: if $f\in\mathcal{H}_k\setminus\mathcal{H}_{k-1}$
then $f\notin\mathrm{EML}_{k-1}$ (since $\mathrm{EML}_{k-1}\subseteq\mathcal{H}_{k-1}$),
so $\mathrm{depth}(f)\geq k$.
\end{proof}

% =============================================================================
\section{Cross-Reference: Theorem Catalog Labels}

For cross-reference with the monogate theorem catalog:

\begin{description}
  \item[T01] Infinite Zeros Barrier: $\sin(x)$ has no finite real EML tree.
  \item[T02] Every elementary function is an exact EML tree (Odrzywołek 2026).
  \item[T07] Euler Gateway: $\sin(x)/\cos(x)=1$ node over $\mathbb{C}$.
  \item[T09] Negation in 2 nodes: $-x=\mathrm{exl}(0,\mathrm{deml}(x,1))$.
  \item[T10] Multiplication in 2 nodes via EXL bridge.
  \item[T11] Addition/subtraction in 3 nodes.
  \item[T16] Tight Zeros Bound: depth-$k$ EML tree has $\leq 2^k$ real zeros.
  \item[T20] $\ln(x)$ construction in 2 EML nodes.
  \item[T21] EDL completeness: $1/x$ in 2 EML nodes.
  \item[T30] EML Depth Spectrum Theorem (this document).
\end{description}

% =============================================================================
\section*{Acknowledgments}

This work builds on Odrzywołek (2026), arXiv:2603.21852, which proved that
EML is exactly complete (T02) and initiated the systematic study of the
EML operator family.  The depth spectrum questions emerged from the monogate
research program on one-operator universal computation.

\bigskip
\noindent\textit{Almaguer, A.R.\ (2026). The EML Depth Spectrum Classification.
monogate research. \url{https://monogate.org}}

\end{document}
